\begin{pastexam}[label={px:nccu-114-dsa-05}]{政大資科所 114 學年度 第 5 題}
（單選；每題 4 分）

(4\%) Given a logical expression represented as an expression tree, along with the Boolean values
assigned to the operands, which algorithm is best suited to evaluate the logical expression?
\begin{enumerate}[label=\Alph*.]
\item Pre-order Traversal
\item In-order Traversal
\item Post-order Traversal
\item Breadth-First Search
\item Depth-First Search
\end{enumerate}
\end{pastexam}

\begin{pastsoln}
答案是 C。

運算式樹（資結書第~7.7~節）的葉節點是運算元，內部節點是運算子（$\wedge$、$\vee$、$\neg$）。 一個
運算子節點要算出值，必須先知道它的子樹（運算元）的值，所以要\textbf{先處理子節點、最後處理自己}，
這正是後序走訪（資結書第~7.3~節）：
\begin{pseudo}[unbreakable]{運算式樹求值（後序）}
\Function{Eval}{$x$}
    \If{$x$ 是葉節點}
        \State \Return $x$ 的布林值
    \EndIf
    \State $a \gets \proc{Eval}(x.\mathit{left})$ \Comment{$\neg$ 只有一個子節點時只算這一個}
    \State $b \gets \proc{Eval}(x.\mathit{right})$
    \State \Return $a \mathbin{x.\mathit{op}} b$
\EndFunction
\end{pseudo}
每個節點處理一次，時間 $O(n)$。

逐項檢查：
\begin{itemize}
\item A.~前序先處理運算子，這時兩邊的值都還不知道，不能直接算（印出前序得到的是前序式，要另外
用堆疊才能求值），錯。
\item B.~中序得到的是中序式（而且要補括號），本身不能算出值，錯。
\item C.~後序：子樹算完才處理運算子，對。
\item D.~BFS 一層一層由上往下，處理根的時候子樹都還沒算，錯。
\item E.~DFS 是一大類走訪，前序、中序、後序都是 DFS；題目要的是「最適合」的那一種，就是後序，
E 不夠精確，不選。
\end{itemize}
\end{pastsoln}
